\documentclass[rutwik_proposal.tex]{subfiles}

\begin{document}

\subsection{Social Identity Theory}\label{subsec:socidentity}

In response to the growing attention to theories of indvidualism in social psychology in the 1960s and 70s, there was a movement to address the problems frequently observed with these theories. Many scholars provided compelling evidence that ignoring the role of society in defining a person's identity was the root cause for most of the observed problems \cite{Brewer79, Cancian75}. In this movement, there was a marked shift from assuming a Hobbesian state of nature to assuming a more Humesian state (the Hobbesian state of nature refers to Thomas Hobbes' conception of humans as being completely individualistic and wholly utilitarian. In the absence of law, people are constantly at war with each other since everyone is only trying to further their own individual goals. David Hume had argued in his \textit{A Treatise of Human Nature} that this conception was at best a philosophical fiction and that thinking of humans as anything but social animals was completely misguided \cite{Hume1739}).

One of the major theories that came out of this movement was the social identity theory. The social identity theory was formulated by Henri Tajfel and extended on by his student John Turner. The corner stone of this theory is the assumption of a social identity in addition to a personal identity. Social identity refers to the part of our self-concept that arises out of our identification with one or more groups \cite{Tajfel81, Turner86}. Personal identity on the other hand refers to personal definitions of self, such as character traits and abilities. 

Social identity can be very fluid. We might identify as members of different groups in different situations or at different times and so, our perceptions of our social identity will change correspondingly. It is important to point out that while social and personal identity are defined as being mutually exclusive, they constantly interact with each other. We might, for example, be more inclined to see ourselves as members of groups whose ideology is similar to our own. Conversely, membership to a group might alter our personal attitudes to be more in line with the group's norms. In this sense, the social identity theory is capable of including the socialized actor theory and thus portrays a much more complete picture.

In this theory, social norms are defined in terms of the norms of groups. Norms are the result of an implicit collective consensus about which behaviors are appropriate and which are not. Members of a group conform to the norms of the group in a particular situation inasmuch as they perceive themselves to be a part of the group in that given situation. For example, as a graduate student in the ecology and evolutionary biology department at Princeton, my frame of reference is other graduate students in the department. My social identity is defined in terms of my perception of how similar I am to the other students. As a member of the ecology and evolutionary biology department as a whole, my frame of reference includes my professors, post docs in the department, and staff members. Accordingly, my social identity changes as well. Depending on which of these identities is invoked at any given time, I might be induced to behave in different ways. Figure \ref{fig:soc_id} depicts a simple scheme of how a person could be a member of many different groups, each larger group subsuming all of the smaller ones. It is also possible to belong to groups that do not completely overlap. You could be a supporter of two different football teams with each fan base having its own set of norms and shared beliefs.

\begin{figure}[h]
\centering
\includegraphics{Images/multi_id}
\caption{Personal and social identities \cite{Brewer91}}
\label{fig:soc_id}
\end{figure}

Identification with a group leads to depersonalization. When we see ourselves as members of a group, we tend to focus on how similar we are to other members of the group. Group identification also serves to differentiate us from members of other groups. Conversely, our personal identities are comprised of perceived differences between us and members of groups to which we belong \cite{Brewer91}. Group membership can thus be made salient either through emphasizing similarities to the in-group (group that we are members of) or by bringing up differences from the out-group (group that we are not members of). Similarly, personal identity can be made salient by stressing differences from members of the in-group. These mechanisms have proved to be fertile ground for experimentation.

The social identity model has enjoyed much empirical support. In the Stanford Prison Experiment, one of the most famous and morally dubious psychology experiments of the 20th century, Zimbardo showed that perfectly normal college students could be induced to behave in inexplicably terrible ways by simply assigning them different group memberships \cite{Zimbardo07}. Marilynn Brewer, through numerous lab experiments, has also shown the importance of the perception of group membership. Depending on whether social or personal identity is made salient, subjects in lab experiments choose to either support group goals, such as the sustainable use of public goods, or not \cite{Brewer91, Kramer84, Brewer79}.

Although this theory has considerable explanatory power, there are situations in which it does fall short. There are instances, particularly in lab experiments, where there is not much room for identification with a group. However, cooperation is still observed in these situations if the participants are allowed to converse amongst themselves prior to the experiment \cite{Dawes00}. It seems then, that perhaps group identity, at least in cases where group memberships are at best weak, is not as important as getting more information about others. Maybe group norms take on more salience when people either do not have much information about the outcomes of the choices they have to make or if individuals do not trust their personal information. It seems obvious that group norms should hold more sway over those who have stronger ties to the group, but other than a few observations in Schanck's seminal paper on the different communities in Elm Hollow \cite{Schanck32}, I have not found direct evidence to support this hypothesis. Additionally, the theory offers no clues as to how we make decisions in cases where we have to choose between our allegiance to different groups.

Another issue with the social identity theory is that it does not explain why people want to belong to groups. While some group memberships are forced upon us (nationality, race, color), we have considerable choice in relation to other groups (democrats vs. republicans, which sports team to support, whether to be a vegetarian or not). Do we choose groups based on the social or other benefits we foresee enjoying in lieu of membership? While this might explain some of our choices, it certainly does not explain why many of us choose to embrace marginal group membership where there are no tangible benefits. It seems that group choice could depend significantly on factors that are difficult to put into rational terms. From a conservation perspective, this is an important aspect to explore. If we want to create groups in which both people and other species are accorded membership, we need to know how to get people to subscribe to these groups. How should we define these groups in order to make them appealing to people?

Finally, while this theory allows for rapid norm changes (if all or most members of a group adopt a new norm, the group norm should change as well), it does not reveal anything about which kinds of norms are more or less susceptible to change. Why do fads, memes, and fashion trends change so quickly while even inefficient norms such as female genital mutilation linger despite extensive efforts to change them \cite{Mackie96, Shell-Duncan00}?

Notwithstanding these and other shortcomings, the social identity theory has some important implications for using the power of social norms to do more effective conservation. For example, this theory tells us that whatever policies we design have to make group membership salient. Focusing on individual level moral values as opposed to group ethics might not be a very effective approach. At the same time, the groups that these interventions are aimed to address cannot be very big and they should be precisely defined. Research has shown that positive evaluation of a group is inversely correlated with how large the group is proportionate to the out group. Additionally, evaluation is directly correlated with how sharp the distinction between the in group and out group is \cite{Mullen92}. More positive evaluations of a group should translate into increased conformity with group norms. So, appealing to people as members of a majority will most likely not be as effective as a multi-pronged approach that addresses different small, well defined groups separately.

What this theory does not help with is how norms within a group change. It tells us that norms can change and are much more flexible than in the Parsonian view, but it sheds no light on mechanisms through which change occurs. Do groups have a nucleus in the form of group leaders that drive change? Is change a consensus arrived at via democratic processes? Work done on norm interventions are more informative about these questions than are the different theories of norms themselves. However, there are no definitive answers and the answers appear to be situation specific. In order to design appropriate conservation related norm interventions, it is imperative that we get a better handle on the answers to these and related questions. I hope to start contributing to this knowledge during my time as a PhD student. 

\end{document}